\documentclass[11pt,a4paper]{article}
\usepackage[margin=2.5cm]{geometry}
\usepackage{amsmath,amssymb,amsthm}
\usepackage{booktabs}
\usepackage{microtype}
\usepackage{xcolor}
\definecolor{cA}{HTML}{1F4E8C}
\definecolor{cB}{HTML}{B3261E}
\newcommand{\seq}[1]{\texttt{\small #1}}
\newcommand{\C}{\mathcal{C}}
\newcommand{\Cp}{\mathcal{C}^{+}}
\usepackage{parskip}

\theoremstyle{definition}
\newtheorem{definition}{Definition}
\theoremstyle{plain}
\newtheorem{theorem}[definition]{Theorem}
\newtheorem{lemma}[definition]{Lemma}
\newtheorem{corollary}[definition]{Corollary}
\newtheorem{proposition}[definition]{Proposition}

\title{Sturmian sequences of $\sqrt n$: the proofs}
\author{Robin Langer}
\date{}

\begin{document}
\maketitle
\vspace{-2em}

This is the proof-dense companion to \texttt{sturmian.tex}. Every claim made
there is proved here. Two claims that \texttt{sturmian.tex} states for
$n = 2$ and $n = 7$ alone are proved for every $n$ in \S9, and one statement
there is false as written; \S6 gives the counterexample and the repair.

\section{Codings}

Fix an irrational $\alpha \in (0,1)$. Write $\{y\}$ for $y - \lfloor y \rfloor$
and, for $x \in [0,1)$,
\[
  x_k \;=\; \{x + k\alpha\} , \qquad k \geq 0 .
\]

\begin{definition}
The \emph{lower} and \emph{upper codings} of $x$ are the sequences
\[
  \C(\alpha,x)_k \;=\;
  \begin{cases} 1 & x_k \in [1-\alpha,\,1) \\ 0 & \text{otherwise,}\end{cases}
  \qquad\qquad
  \Cp(\alpha,x)_k \;=\;
  \begin{cases} 1 & x_k \in (1-\alpha,\,1) \text{ or } x_k = 0 \\ 0 & \text{otherwise.}\end{cases}
\]
\end{definition}

The two codings read the same circle with the same rotation and differ only in
which endpoint of the arc of length $\alpha$ carries the letter $1$.

\begin{lemma}\label{lem:mech}
For $x \in [0,1)$,
\[
  \C(\alpha,x)_k \;=\; \lfloor (k+1)\alpha + x \rfloor - \lfloor k\alpha + x \rfloor ,
  \qquad
  \Cp(\alpha,x)_k \;=\; \lceil (k+1)\alpha + x \rceil - \lceil k\alpha + x \rceil .
\]
\end{lemma}

\begin{proof}
Put $m = \lfloor k\alpha + x\rfloor$ and $f = x_k$, so $k\alpha + x = m + f$ and
$(k+1)\alpha + x = m + f + \alpha$ with $f + \alpha \in [0,2)$. Hence
$\lfloor (k+1)\alpha+x\rfloor - m = \lfloor f + \alpha \rfloor$, which is $1$
exactly when $f \geq 1-\alpha$.

For the second, $\lceil m + f \rceil = m$ when $f = 0$ and $m+1$ otherwise, so
$\lceil m+f+\alpha\rceil - \lceil m+f\rceil = 1$ exactly when $f = 0$ or
$f + \alpha > 1$.
\end{proof}

So $\C(\alpha,x)$ is the sequence of \texttt{sturmian.tex} \S2 with $\rho = x$.

\begin{lemma}\label{lem:lowup}
$\C(\alpha,x)$ and $\Cp(\alpha,x)$ are equal unless $x_{k_0} = 0$ for some
$k_0 \geq 0$, which happens for at most one $k_0$. In that case they differ
exactly at the positions $k_0 - 1$ and $k_0$, where they read \seq{10} and
\seq{01}.
\end{lemma}

\begin{proof}
The letter rules differ only at $x_k \in \{0,\, 1-\alpha\}$. If $x_k = 1-\alpha$
then $x_{k+1} = 0$, and if $x_{k_0} = 0$ with $k_0 \geq 1$ then
$x_{k_0 - 1} = 1 - \alpha$. Two indices with $x_k = 0$ would give
$(j-k)\alpha \in \mathbb{Z}$.
\end{proof}

\begin{lemma}\label{lem:count}
For $x \in [0,1)$ and $N \geq 1$ the number of $1$s among the first $N$ letters
of $\C(\alpha,x)$ is $\lfloor N\alpha + x\rfloor$.
\end{lemma}

\begin{proof}
Sum Lemma~\ref{lem:mech} over $k < N$ and use $\lfloor x \rfloor = 0$.
\end{proof}

\begin{corollary}\label{cor:slope}
The density of $1$s in $\C(\alpha,x)$ is $\alpha$. The same holds for
$\Cp(\alpha,x)$.
\end{corollary}

\begin{proof}
$\lfloor N\alpha + x\rfloor / N \to \alpha$. The upper coding differs in at
most two positions.
\end{proof}

The density is called the \emph{slope}. It is a property of the sequence alone,
so two codings that are equal as sequences have equal slopes. Every use of
``solve for $\alpha$'' below rests on this corollary.

\begin{lemma}\label{lem:shift}
$\C(\alpha,x)_{k+1} = \C(\alpha,\{x+\alpha\})_k$, and likewise for $\Cp$.
\end{lemma}

\begin{proof}
$\{x + (k+1)\alpha\} = \{\{x+\alpha\} + k\alpha\}$.
\end{proof}

\section{The orbit is dense}

\begin{lemma}\label{lem:dense}
For irrational $\alpha$ and any $x \in [0,1)$, the set $\{x_k : k \geq 0\}$ is
dense in $[0,1)$.
\end{lemma}

\begin{proof}
Let $\varepsilon > 0$. By Dirichlet there are integers $q \geq 1$ and $p$ with
$|q\alpha - p| < 1/q < \varepsilon$; put $\delta = |q\alpha - p| > 0$. The
points $x_0,\, x_q,\, x_{2q},\, \ldots$ differ consecutively by $\pm\delta$
modulo $1$, with a fixed sign, so the first $\lceil 1/\delta\rceil + 1$ of them
meet every interval of length $\delta$. Take $\delta < \varepsilon$.
\end{proof}

\begin{corollary}\label{cor:inj}
$x \mapsto \C(\alpha,x)$ is injective on $[0,1)$.
\end{corollary}

\begin{proof}
Let $x < y$ and suppose the codings agree. By Lemma~\ref{lem:count},
$\lfloor N\alpha + x\rfloor = \lfloor N\alpha+y\rfloor$ for every $N \geq 1$.
Choose $N$ with $\{N\alpha + x\} > 1 - (y-x)$; then
$N\alpha + x < \lceil N\alpha+x\rceil \leq N\alpha + y$.
\end{proof}

\section{Complexity}

\begin{definition}
A \emph{factor} of length $m$ of a sequence $\ell$ is a block
$\ell_k \ell_{k+1}\cdots \ell_{k+m-1}$. Let $F_\ell(m)$ be the number of
distinct factors of length $m$. The sequence is \emph{Sturmian} when
$F_\ell(m) = m+1$ for every $m \geq 1$.
\end{definition}

\begin{theorem}\label{thm:complexity}
$F_{\C(\alpha,x)}(m) = m+1$ for every $m \geq 1$.
\end{theorem}

\begin{proof}
For $j \geq 1$ and $t \in [0,1)$,
\[
  \lfloor t + j\alpha \rfloor \;=\; \lfloor j\alpha \rfloor +
  \begin{cases} 1 & t \geq \{-j\alpha\} \\ 0 & t < \{-j\alpha\} ,\end{cases}
\]
because $t + j\alpha = \lfloor j\alpha\rfloor + (\{j\alpha\} + t)$ with
$\{j\alpha\} + t \in [0,2)$, and $\{j\alpha\} + t \geq 1$ exactly when
$t \geq 1 - \{j\alpha\} = \{-j\alpha\}$.

Let $w$ be the factor of length $m$ at position $k$. By Lemma~\ref{lem:shift}
its letters are $w_j = \C(\alpha,x_k)_j$ for $j < m$, so by
Lemma~\ref{lem:count} the partial sums of $w$ give $\lfloor x_k + j\alpha\rfloor$
for every $j \leq m$. By the display, $w$ therefore determines, for each
$j \in \{1,\ldots,m\}$, whether $x_k \geq \{-j\alpha\}$; and conversely those
$m$ comparisons determine $w$.

The $m$ points $\{-j\alpha\}$, $1 \leq j \leq m$, are distinct and lie in
$(0,1)$, so they cut $[0,1)$ into $m+1$ intervals of positive length. Factors
of length $m$ are in bijection with those intervals that contain some $x_k$, and
by Lemma~\ref{lem:dense} every one of them does.
\end{proof}

\begin{corollary}\label{cor:complexityup}
$F_{\Cp(\alpha,x)}(m) = m+1$ for every $m \geq 1$.
\end{corollary}

\begin{proof}
One can rerun the proof above with $\lceil\cdot\rceil$ for
$\lfloor\cdot\rfloor$: the comparisons become $x_k > \{-j\alpha\}$ and the
intervals are half-open at the other end. But that is more work than is needed,
because an upper coding is the bitwise complement of a lower coding:
\[
  \Cp(\alpha,x)_k + \C(1-\alpha,-x)_k = 1 \qquad (k \geq 0).
\]
Indeed $\Cp(\alpha,x)_k = 1$ exactly when
$\{x+k\alpha\} \in (1-\alpha,1) \cup \{0\}$, whose complement in $[0,1)$ is
$(0,1-\alpha]$; and $\{-x + k(1-\alpha)\} = \{-(x+k\alpha)\}$, so the
reflection $v \mapsto 1-v$ sends the orbit of $x$ at slope $\alpha$ to the
orbit of $-x$ at slope $1-\alpha$, carrying $(0,1-\alpha]$ onto
$[\alpha,1) = [1-(1-\alpha),1)$ --- the arc of the lower coding of slope
$1-\alpha$. The exceptional point $v = 0$ is fixed by the reflection and lies
outside $[\alpha,1)$, as it must. Complementation is a bijection on words, and
$1-\alpha$ is again irrational and in $(0,1)$, so
Theorem~\ref{thm:complexity} gives the count.
\end{proof}

The second proof is the one formalised. In \texttt{lean/Sturmian.lean} the
displayed identity is
\texttt{codeUp\_add\_code\_reflect}, and the corollary is
\texttt{ncard\_factors\_codeUp}. The first proof would have cost a parallel
development: the natural range of the base point for
an upper coding is $(0,1]$ rather than $[0,1)$, since $\Cp$ reads
$\{x+k\alpha\} = 0$ as the letter $1$.

For the Fibonacci sequence of \texttt{sturmian.tex} \S1 and $m = 3$ the four
intervals are cut by $\{-\alpha\}, \{-2\alpha\}, \{-3\alpha\}$ with
$\alpha = 1/\varphi$, and the four factors are \seq{101}, \seq{011}, \seq{110},
\seq{010}. That is the table in \S1.

\section{Gaps}

Write
\[
  \frac{1}{\alpha} \;=\; a_1 + \beta, \qquad a_1 = \lfloor 1/\alpha\rfloor \geq 1,
  \qquad \beta = \{1/\alpha\} \in (0,1) .
\]
Then $a_1 \alpha < 1 < (a_1+1)\alpha$, which is the inequality of
\texttt{sturmian.tex} \S3.

Let $I = [1-\alpha, 1)$, the arc carrying the letter $1$, and coordinatise it by
\[
  x \;=\; (1-\alpha) + \alpha u , \qquad u \in [0,1) .
\]

\begin{theorem}\label{thm:return}
Let $x \in I$ with coordinate $u$. The least $t \geq 1$ with
$\{x + t\alpha\} \in I$ is
\[
  t \;=\; \begin{cases} a_1 + 1 & u < \beta \\ a_1 & u \geq \beta, \end{cases}
\]
and the coordinate of $\{x+t\alpha\}$ is $\{u + 1 - \beta\}$.
\end{theorem}

\begin{proof}
$x + \alpha \in [1, 1+\alpha)$, so $\{x+\alpha\} = x + \alpha - 1 = \alpha u$,
and for $s \geq 1$,
\[
  \{x + s\alpha\} \;=\; \alpha(u + s - 1)
  \qquad\text{as long as } \alpha(u+s-1) < 1 .
\]
Now $\alpha(u+s-1) < 1 \iff s - 1 < 1/\alpha - u = a_1 + \beta - u$, and
$\alpha(u+s-1) \geq 1-\alpha \iff u + s \geq a_1 + \beta$. So
$\{x+s\alpha\} \in I$ exactly when
\[
  a_1 + \beta - u \;\leq\; s \;<\; a_1 + \beta - u + 1 ,
\]
whose unique integer solution is $s = \lceil a_1 + \beta - u\rceil$. For
$u < \beta$ that is $a_1+1$ and for $u \geq \beta$ it is $a_1$; in both cases
$s$ is the least such $t$ because the displayed formula for
$\{x+s\alpha\}$ is valid and below $1-\alpha$ for all smaller $s$.

The new coordinate is $(\alpha(u+t-1) - (1-\alpha))/\alpha = u + t - 1/\alpha$,
which is $u+1-\beta$ when $t = a_1+1$ and $u - \beta$ when $t = a_1$.
\end{proof}

\begin{corollary}\label{cor:gaps}
In $\C(\alpha,x)$ the gap between consecutive $1$s is $a_1$ or $a_1+1$, so the
number of $0$s between them is $a_1 - 1$ or $a_1$.
\end{corollary}

\begin{definition}
Let $k_0 < k_1 < k_2 < \cdots$ be the positions of the $1$s of $\ell$. The
\emph{derived sequence} $\ell'$ is
\[
  \ell'_i \;=\; \begin{cases}
    1 & \text{$a_1$ zeros lie between the $i$th and $(i{+}1)$th $1$} \\
    0 & \text{$a_1 - 1$ do.}
  \end{cases}
\]
\end{definition}

By Theorem~\ref{thm:return}, $\ell'_i = 1$ exactly when the coordinate $u_i$ of
$x_{k_i}$ satisfies $u_i < \beta$, and
\[
  u_{i+1} \;=\; \{u_i + 1 - \beta\} .
\]

\begin{lemma}[block decomposition]\label{lem:blocks}
Define $\tau_a(1) = \seq{1}\,\seq{0}^a$ and $\tau_a(0) = \seq{1}\,\seq{0}^{a-1}$.
Then
\[
  \ell \;=\; \seq{0}^{k_0} \cdot \tau_{a_1}(\ell') .
\]
\end{lemma}

\begin{proof}
The $i$th block of $\ell$ starting at $k_i$ is a $1$ followed by the zeros
before the next $1$, of which there are $a_1$ if $\ell'_i = 1$ and $a_1-1$ if
$\ell'_i = 0$.
\end{proof}

Deriving and $\tau_{a_1}$ are inverse to each other up to the prefix
$\seq{0}^{k_0}$. That prefix is the only reason a shift appears anywhere below.

\section{The derivation theorem}

\begin{theorem}\label{thm:derive}
Let $u_0$ be the coordinate of $x_{k_0}$, the first $1$. Then
\[
  \C(\alpha,x)' \;=\; \Cp(\beta,\, \{-u_0\}) ,
  \qquad\qquad
  \Cp(\alpha,x)' \;=\; \C(\beta,\, 1 - u_0) .
\]
\end{theorem}

\begin{proof}
For the lower coding put $v_i = \{-u_i\}$. By Theorem~\ref{thm:return},
$v_{i+1} = \{v_i + \beta\}$, so $v_i = \{v_0 + i\beta\}$. And
\[
  \ell'_i = 1 \iff u_i \in [0,\beta)
  \iff -u_i \in (-\beta, 0]
  \iff v_i \in (1-\beta,1) \text{ or } v_i = 0 ,
\]
which is the letter rule of $\Cp(\beta, v_0)$.

For the upper coding the arc is $I^+ = (1-\alpha,1]$, coordinatised by
$u \in (0,1]$ with $u = 1$ meaning $x \equiv 0$. The computation of
Theorem~\ref{thm:return} is unchanged except that $\{x+s\alpha\} \in I^+$ reads
$u + s > a_1 + \beta$ and $u + s - 1 \leq a_1 + \beta$, so
$t = a_1 + 1$ exactly when $u \leq \beta$. Put $v_i = 1 - u_i \in [0,1)$; then
$v_{i+1} = \{v_i + \beta\}$ and
\[
  \ell'_i = 1 \iff u_i \leq \beta \iff v_i \geq 1 - \beta ,
\]
which is the letter rule of $\C(\beta, v_0)$.
\end{proof}

Deriving sends the slope $\alpha$ to $\beta = \{1/\alpha\}$ and exchanges the
two codings. The exchange is invisible for all but countably many $x$
(Lemma~\ref{lem:lowup}), and it is exactly what makes the statement of
Theorem~\ref{thm:derive} an equality rather than an equality off a small set.

\begin{corollary}\label{cor:derivesturmian}
The derived sequence of a coding is a coding, hence Sturmian, and its slope is
$\{1/\alpha\}$.
\end{corollary}

\begin{corollary}\label{cor:gauss}
The slopes of $\ell, \ell', \ell'', \ldots$ are the orbit of $\alpha$ under the
Gauss map $\alpha \mapsto \{1/\alpha\}$, so the block widths $a_1$ encountered
are the partial quotients of $\alpha = [\,0;\, a_1, a_2, a_3, \ldots\,]$.
\end{corollary}

\section{Morse and Hedlund}

One direction is Theorem~\ref{thm:complexity}: every coding with irrational
slope is Sturmian. The statement in \texttt{sturmian.tex} \S2 asserts the
converse in the form
\[
  \ell_k \;=\; \lfloor (k+1)\alpha + \rho\rfloor - \lfloor k\alpha + \rho\rfloor,
  \qquad \rho \in [0,1) ,
\]
that is, $\ell = \C(\alpha,\rho)$. In that form it is false.

Let $\alpha = 1/\varphi = (\sqrt 5 - 1)/2$ and let $c = \C(\alpha,\alpha)$, the
Fibonacci sequence. Let $\seq{1}c$ be $c$ with a $1$ prefixed.

\begin{theorem}\label{thm:counterex}
$\seq{1}c = \Cp(\alpha,0)$. It is Sturmian, and there is no $x \in [0,1)$ with
$\seq{1}c = \C(\alpha', x)$ for any irrational $\alpha' \in (0,1)$.
\end{theorem}

\begin{proof}
$\Cp(\alpha,0)_0 = 1$ because $x_0 = 0$, and
$\Cp(\alpha,0)_{k+1} = \Cp(\alpha,\alpha)_k = \C(\alpha,\alpha)_k$ by
Lemma~\ref{lem:shift} and Lemma~\ref{lem:lowup}, the orbit of $\alpha$ missing
$0$. So $\seq{1}c = \Cp(\alpha,0)$, and it is Sturmian by
Corollary~\ref{cor:complexityup}.

By Lemma~\ref{lem:count} the first $N-1$ letters of $c$ carry
$\lfloor (N-1)\alpha + \alpha\rfloor = \lfloor N\alpha\rfloor$ ones, so the first
$N$ letters of $\seq{1}c$ carry $1 + \lfloor N\alpha\rfloor$.

Suppose $\seq{1}c = \C(\alpha',x)$. The slopes agree by
Corollary~\ref{cor:slope}, so $\alpha' = \alpha$, and Lemma~\ref{lem:count}
gives $\lfloor N\alpha + x\rfloor = 1 + \lfloor N\alpha \rfloor$ for every
$N \geq 1$. Since $x \in [0,1)$ this forces $\{N\alpha\} + x \geq 1$ for every
$N \geq 1$. By Lemma~\ref{lem:dense} there is $N$ with $\{N\alpha\} < 1-x$.
\end{proof}

The repair is to admit both codings.

\begin{theorem}[Morse and Hedlund]\label{thm:mh}
A sequence is Sturmian if and only if it is $\C(\alpha,x)$ or $\Cp(\alpha,x)$
for some irrational $\alpha \in (0,1)$ and some $x \in [0,1)$.
\end{theorem}

The forward implication is Theorem~\ref{thm:complexity} and
Corollary~\ref{cor:complexityup}. For the converse, run \S4 and \S5 backwards.
Steps 1 to 4 are below; step 5 is \S7.

\begin{enumerate}
\item $F_\ell(1) = 2$, so both letters occur. $F_\ell(2) = 3$, so exactly one
  word of length $2$ is absent, and it is $\seq{00}$ or $\seq{11}$: were
  $\seq{01}$ absent, every $1$ would precede every $0$ and
  $\ell = \seq{1}^p\seq{0}^\infty$; were $\seq{10}$ absent,
  $\ell = \seq{0}^p\seq{1}^\infty$. Either has complexity $p+1$ for all
  $m > p$. \emph{Proved.}
\item The lengths of the maximal blocks of $0$s take exactly two values, and
  they are consecutive, say $a-1$ and $a$. \emph{Proposition~\ref{prop:blocks}
  below.}
\item Given step 2, Lemma~\ref{lem:blocks} applies verbatim: $\ell$ factors as
  $\seq{0}^{k_0}\tau_a(\ell')$ with $\ell'$ a sequence over
  $\{\seq{0},\seq{1}\}$. \emph{Proved.}
\item $\ell'$ is again Sturmian. \emph{Proposition~\ref{prop:derived} below};
  that both of its letters occur is the case $m = 1$.
\item Iterating pins a single $\alpha$ and a single $x$. \emph{\S7:}
  Lemma~\ref{lem:slope} and Theorem~\ref{thm:step5}.
\end{enumerate}

Step 1 delivers $\seq{00}$ or $\seq{11}$, and step 2 is stated for $\seq{11}$.
Exchanging the letters costs nothing: by the identity in the proof of
Corollary~\ref{cor:complexityup}, the complement of $\C(\beta,y)$ is
$\Cp(1-\beta,-y)$, and Theorem~\ref{thm:mh} admits both codings.

Step 2 is a counting argument. Call a factor $u$ of $\ell$ \emph{right-special}
when $u\seq{0}$ and $u\seq{1}$ are both factors.

\begin{lemma}\label{lem:rs}
$F_\ell(m+1) - F_\ell(m)$ is the number of right-special factors of length $m$.
\end{lemma}

\begin{proof}
Every factor of length $m+1$ is $ua$ for a factor $u$ of length $m$ and a letter
$a$, and every factor of length $m$ has at least one right extension: it occurs
at some $k$, and $\ell_{k+m}$ extends it. So $ua \mapsto u$ is onto the factors
of length $m$, with two elements over $u$ exactly when $u$ is right-special.
\end{proof}

\begin{corollary}\label{cor:onesp}
$F_\ell(m) = m+1$ for every $m$ if and only if $\ell$ has exactly one
right-special factor of each length.
\end{corollary}

The pair that does the work is $\seq{0}^{j+1}$ and $\seq{1}\seq{0}^{j}$. They
have the same length and differ in their first letter, so at most one of them is
right-special.

\begin{proposition}\label{prop:blocks}
Let $\ell$ be Sturmian with $\seq{11}$ not a factor, and let $a$ be the length of
the longest block of $\seq{0}$s. Then $a \geq 2$, the blocks of $\seq{0}$s
between consecutive $\seq{1}$s have lengths $a-1$ and $a$ and both occur, and the
block before the first $\seq{1}$ has length at most $a$.
\end{proposition}

\begin{proof}
$\ell$ has infinitely many $\seq{1}$s: otherwise it ends in $\seq{0}^\infty$ and
$F_\ell$ is bounded. Let
\[
  S \;=\; \{\, j \;:\; \seq{1}\seq{0}^{j}\seq{1} \text{ is a factor of } \ell \,\},
\]
the set of interior block lengths. $S$ is nonempty, and $j \geq 1$ for $j \in S$
because $\seq{11}$ is absent.

The blocks of $\seq{0}$s are bounded. Were they not, $S$ would be unbounded,
only one block being initial. Take $j \in S$ and then $s \in S$ with
$s \geq j+2$. Now $\seq{1}\seq{0}^{j}\seq{1}$ and $\seq{1}\seq{0}^{j+1}$ are
factors, so $\seq{1}\seq{0}^{j}$ is right-special; and $\seq{0}^{j+2}$ and
$\seq{0}^{j+1}\seq{1}$ are factors, so $\seq{0}^{j+1}$ is right-special. Two of
them, of length $j+1$, against Corollary~\ref{cor:onesp}.

So $a$ is defined. Every block is finite, so a block of length $\geq j$ is
followed by a $\seq{1}$: hence $\seq{0}^{j}$ is a factor exactly for
$j \leq a$, and $\seq{0}^{j}\seq{1}$ is a factor for $1 \leq j \leq a$.

Let $j \in S$ with $j < \max S$. Then $\seq{1}\seq{0}^{j}\seq{1}$ and
$\seq{1}\seq{0}^{j+1}$ are factors, so $\seq{1}\seq{0}^{j}$ is right-special. If
$j \leq a-2$ then $\seq{0}^{j+2}$ and $\seq{0}^{j+1}\seq{1}$ are factors, so
$\seq{0}^{j+1}$ is right-special as well, and Corollary~\ref{cor:onesp} is
contradicted. So
\[
  j \in S,\ j < \max S \;\Longrightarrow\; j \geq a-1 .
\]

If $\max S < a$, every $j \in S$ below $\max S$ would satisfy
$j \geq a-1 \geq \max S$, so there is none and $S = \{\max S\}$. A single block
length makes $\ell$ eventually periodic, which bounds $F_\ell$. So $\max S = a$,
and the display gives $S \subseteq \{a-1,a\}$. Again $S$ is not a single value,
so $S = \{a-1,a\}$, and $a-1 \geq 1$.

The initial block has length at most $a$ by the definition of $a$.
\end{proof}

Proposition~\ref{prop:blocks} is weaker than balancedness and is what steps 3
to 5 use. The classical route proves balancedness in full, by a
minimal-counterexample induction producing a palindrome $w$ with
$\seq{0}w\seq{0}$ and $\seq{1}w\seq{1}$ both factors; that is not needed here.

Step 4 needs three more facts, all about counting, and then one bridge.

\begin{lemma}\label{lem:evper}
An eventually periodic sequence has bounded factor complexity.
\end{lemma}

\begin{proof}
With preperiod $q$ and period $p$, every factor is a factor at some position
below $q+p$, so $F_\ell(n) \leq q+p$.
\end{proof}

\begin{lemma}\label{lem:aper}
If $\ell$ has no right-special factor of length $m$, it is eventually periodic.
\end{lemma}

\begin{proof}
Every factor of length $m$ has exactly one right extension, so the window
$u_k = \ell_k\cdots\ell_{k+m-1}$ determines $u_{k+1}$. The windows lie in a
finite set, so $u_0, u_1, \ldots$ is eventually periodic, and $\ell$ is read off
its windows.
\end{proof}

\begin{lemma}\label{lem:rschain}
If $\ell$ has exactly one right-special factor of each length, each is a suffix
of the next.
\end{lemma}

\begin{proof}
The suffix of length $m$ of a right-special factor of length $m+1$ is
right-special.
\end{proof}

\begin{lemma}\label{lem:bridge}
Let $\ell$, $a$ and $\ell'$ be as in step 3. Then $w$ is a right-special factor
of $\ell'$ if and only if $\tau_a(w)\seq{1}\seq{0}^{a-1}$ is a right-special
factor of $\ell$.
\end{lemma}

\begin{proof}
The $\seq{1}$s of $\ell$ are exactly the block starts, and $\tau_a(w)$ begins
with $\seq{1}$, so every occurrence of $\tau_a(w)$ in $\ell$ begins at a block
start. So $w\seq{0}$ is a factor of $\ell'$ exactly when
$\tau_a(w)\seq{1}\seq{0}^{a-1}\seq{1}$ is a factor of $\ell$, and $w\seq{1}$
exactly when $\tau_a(w)\seq{1}\seq{0}^{a}$ is one, no block being longer
than $a$.
\end{proof}

\begin{proposition}\label{prop:derived}
Let $\ell$, $a$ and $\ell'$ be as in step 3. Then $\ell'$ is Sturmian.
\end{proposition}

\begin{proof}
At most one right-special factor of each length. Let $w$ and $w'$ be
right-special factors of $\ell'$ of length $m$. By Lemma~\ref{lem:bridge} both
$\tau_a(w)\seq{1}\seq{0}^{a-1}$ and $\tau_a(w')\seq{1}\seq{0}^{a-1}$ are
right-special factors of $\ell$, so by Lemma~\ref{lem:rschain} one is a suffix
of the other, and after dropping the common $\seq{1}\seq{0}^{a-1}$ so is one of
$\tau_a(w)$, $\tau_a(w')$. A suffix of $\tau_a(w')$ beginning with $\seq{1}$
begins at a block start, so it is $\tau_a(v)$ for a suffix $v$ of $w'$; counting
blocks, $v = w'$. So $\tau_a(w) = \tau_a(w')$ and $w = w'$.

At least one right-special factor of each length. Were there none of length $m$,
$\ell'$ would be eventually periodic by Lemma~\ref{lem:aper}, hence so would
$\ell = \seq{0}^{k_0}\tau_a(\ell')$, and $F_\ell$ would be bounded by
Lemma~\ref{lem:evper}.

Corollary~\ref{cor:onesp} now gives $F_{\ell'}(m) = m+1$.
\end{proof}

\section{Step 5}

Normalise so that \seq{11} is absent: by step 1 and the identity in the proof of
Corollary~\ref{cor:complexityup}, the complement of $\C(\beta,y)$ is
$\Cp(1-\beta,-y)$, and both codings are admitted. Then
Proposition~\ref{prop:blocks} gives $a \geq 2$, and every level below is
normalised the same way.

Write $\Delta_\alpha(u) = \lvert u\rvert_1 - \alpha\lvert u\rvert$ for a finite
word $u$, where $\lvert u \rvert_1$ counts the $1$s.

\begin{lemma}\label{lem:slope}
Let $a_1, a_2, \ldots \geq 2$ be the block parameters of the successive derived
sequences and $\varepsilon_1, \varepsilon_2, \ldots \in \{0,1\}$ record which
levels were complemented. There is exactly one $\alpha \in (0,1)$ compatible
with this data, and it is irrational.
\end{lemma}

\begin{proof}
A slope and the slope below it are related by $\alpha = 1/(a_1 + \beta)$ and
$\beta \mapsto \beta$ or $1-\beta$. On $[0,1]$ the map $\beta \mapsto
1/(a_1+\beta)$ has derivative $-1/(a_1+\beta)^2$ of modulus at most
$1/a_1^2 \leq 1/4$, and $\beta \mapsto 1 - \beta$ is an isometry. So the first
$j$ levels confine $\alpha$ to an interval of length at most $4^{-j}$.

Were $\alpha = p/q$ in lowest terms, then $\beta = q/p - a_1$ has denominator
$p < q$, and complementing does not change a denominator. Denominators cannot
descend forever, and every level has a block parameter by
Proposition~\ref{prop:derived}.
\end{proof}

\begin{lemma}\label{lem:dev}
$\Delta_\alpha(\tau_a(w)) = -\alpha\,\Delta_\beta(w)$.
\end{lemma}

\begin{proof}
$\lvert \tau_a(w)\rvert = a\lvert w\rvert + \lvert w\rvert_1$ and
$\lvert \tau_a(w)\rvert_1 = \lvert w\rvert$, so
\[
  \Delta_\alpha(\tau_a(w)) = \lvert w\rvert(1 - a\alpha) - \alpha\lvert w\rvert_1
  = \alpha\beta\lvert w\rvert - \alpha\lvert w\rvert_1 .
\]
\end{proof}

\begin{proposition}\label{prop:bound}
$\lvert \Delta_\alpha(u)\rvert < 1$ for every factor $u$ of $\ell$.
\end{proposition}

\begin{proof}
By induction on $\lvert u\rvert$, over all levels at once; $\alpha$ and $\beta$
are the slopes of $\ell$ and $\ell'$, and $1/\alpha = a + \beta$.

Write $u = \seq{0}^s\,\tau_a(u')\,q$ with $u'$ a factor of $\ell'$ and $q$ empty
or $q = \seq{1}\seq{0}^t$ a \emph{proper} prefix of its block. A block has
$a-1$ or $a$ zeros, so $s \leq a$ and $t \leq a-1$.

For $u'$ empty, $\Delta_\alpha(u) = -\alpha s$ or $1 - \alpha(s+1+t)$, and
$\alpha a < 1$ bounds both: $-1 < -\alpha a \leq -\alpha s \leq 0$ and
$-1 < 1 - 2\alpha a \leq 1 - \alpha(s+1+t) \leq 1-\alpha$.

For $u'$ nonempty, $s = a$ forces the block to the left of $u'$ to be the long
one, so $\seq{1}u'$ is a factor of $\ell'$. Also
\[
  \lvert u\rvert \;\geq\; s + a\lvert u'\rvert + \lvert u'\rvert_1
  \;\geq\; a(1 + \lvert u'\rvert) \;>\; \lvert u'\rvert + 1 ,
\]
so the induction hypothesis applies to $u'$ and to $\seq{1}u'$. By
Lemma~\ref{lem:dev},
\[
  \Delta_\alpha(u) \;=\; -\alpha s - \alpha\Delta_\beta(u')
  \;+\; \begin{cases} 1 - \alpha(1+t) & q \neq \epsilon \\ 0 & q = \epsilon .\end{cases}
\]
For the upper bound, $\Delta_\beta(u') > -1$ gives $-\alpha s -
\alpha\Delta_\beta(u') < \alpha$ when $q = \epsilon$, and $s + 1 + t +
\Delta_\beta(u') > 0$ gives $\Delta_\alpha(u) < 1$ when $q \neq \epsilon$.

For the lower bound, $s = a$ gives $\Delta_\beta(\seq{1}u') = 1 - \beta +
\Delta_\beta(u') < 1$, hence $\Delta_\beta(u') < \beta$. Now:
\[
  \begin{array}{llll}
    q = \epsilon, & s \leq a-1: & \alpha(s + \Delta_\beta(u')) < \alpha a & < 1 \\
    q = \epsilon, & s = a: & \alpha(a + \beta) & = 1 \\
    q \neq \epsilon, & s \leq a-1: & \alpha(s+1+t+\Delta_\beta(u')) < 2\alpha a & < 2 \\
    q \neq \epsilon, & s = a: & \alpha(2a + \beta) = \alpha a + 1 & < 2 .
  \end{array}
\]
\end{proof}

\begin{lemma}\label{lem:countup}
The number of $1$s among the first $N$ letters of $\Cp(\alpha,x)$ is
$\lceil N\alpha + x\rceil - \lceil x \rceil$.
\end{lemma}

\begin{proof}
Sum Lemma~\ref{lem:mech} over $k < N$.
\end{proof}

Let $S(N)$ be the number of $1$s among the first $N$ letters of $\ell$, and
\[
  D(N) \;=\; S(N) - N\alpha, \qquad D(0) = 0 .
\]
For $M < N$, $D(N) - D(M) = \Delta_\alpha(u)$ for the factor $u = \ell_M \cdots
\ell_{N-1}$, so Proposition~\ref{prop:bound} says the values of $D$ lie in an
interval of length at most $1$.

\begin{theorem}\label{thm:step5}
Put $x = \sup_{N \geq 0} D(N)$. Then $x \in [0,1]$, and
\[
  \ell \;=\; \begin{cases}
    \C(\alpha, x) & \text{if the supremum is attained} \\
    \Cp(\alpha, x) & \text{if not, and } x < 1 \\
    \Cp(\alpha, 0) & \text{if not, and } x = 1 .
  \end{cases}
\]
\end{theorem}

\begin{proof}
$D(0) = 0$ gives $x \geq 0$, and $D(N) - D(0) < 1$ gives $x \leq 1$.

Attained at $N_0$: then $D(N) \leq x$ and $x - D(N) = D(N_0) - D(N) < 1$ for
every $N$, so $S(N) \leq N\alpha + x < S(N)+1$ and
$\lfloor N\alpha + x\rfloor = S(N)$. Lemma~\ref{lem:mech} gives
$\C(\alpha,x)_k = S(k+1) - S(k) = \ell_k$.

Not attained: then $D(N) < x$ for every $N$, and $x \leq D(N) + 1$. For
$x \in (0,1)$ this reads $S(N) < N\alpha + x \leq S(N)+1$, so
$\lceil N\alpha + x\rceil = S(N) + 1$ and Lemma~\ref{lem:countup} applies with
$\lceil x\rceil = 1$. For $x = 1$ it reads $S(N) \geq N\alpha$ and
$S(N) < N\alpha + 1$, so $\lceil N\alpha \rceil = S(N)$.
\end{proof}

Together with Lemma~\ref{lem:slope} this is step 5, and
Theorem~\ref{thm:mh} is proved. The sequence $\seq{1}c$ of
Theorem~\ref{thm:counterex} is the third case: its $D(N) = 1 - \{N\alpha\}$ has
supremum $1$, attained nowhere.

\section{The substitution}

For $a_1 = 1$ the two blocks are $\tau_1(1) = \seq{10}$ and $\tau_1(0) = \seq{1}$,
so $\tau_1$ is the substitution $\sigma$ of \texttt{sturmian.tex} \S1:
\[
  \sigma(1) = \seq{10}, \qquad \sigma(0) = \seq{1} .
\]

\begin{theorem}\label{thm:fib}
Let $\alpha = 1/\varphi$ and $c = \C(\alpha,\alpha)$. Then $c' = c$,
$\sigma(c) = c$, and $c = \lim_j \sigma^j(\seq{1})$.
\end{theorem}

\begin{proof}
$1/\alpha = \varphi = 1 + 1/\varphi$, so $a_1 = 1$ and $\beta = \alpha$. The
first $1$ of $c$ is at $k_0 = 0$, where $x_0 = \alpha$ has coordinate
$u_0 = (\alpha - (1-\alpha))/\alpha = 2 - 1/\alpha = 1 - \alpha$. So
$\{-u_0\} = \alpha$ and Theorem~\ref{thm:derive} gives
$c' = \Cp(\alpha,\alpha) = \C(\alpha,\alpha) = c$, the middle equality by
Lemma~\ref{lem:lowup}. Lemma~\ref{lem:blocks} with $k_0 = 0$ then reads
$c = \tau_1(c') = \sigma(c)$.

$\sigma(\seq{1}) = \seq{10}$ begins with $\seq{1}$, and if $\sigma^j(\seq{1})$ is a
prefix of $\sigma^{j+1}(\seq{1})$ then applying $\sigma$ keeps it so; hence the
limit exists. It is a prefix of every $\sigma$-fixed sequence beginning with
$\seq{1}$, and $c$ is one.
\end{proof}

Theorem~\ref{thm:fib} is \texttt{fib\_derive}, \texttt{sigma\_pre} and
\texttt{fib\_limit} in \texttt{lean/Sturmian.lean}. Two remarks on the
formalisation. First, $c' = c$ is the instance at the fixed point $\beta =
\alpha$ of the general \texttt{step\_low}, so the only new computation is
$1/\alpha = 1 + \alpha$. Second, $\sigma(c) = c$ and $c = \lim_j
\sigma^j(\seq{1})$ are stated on prefixes --- $\sigma$ applied to the first $m$
letters of $c$ gives the first $k_m$ letters of $c$, where $k_m$ is the position
of the $m$th $1$ --- because $\sigma$ of an infinite sequence is not a function
of position without first inverting the block-start map. Lemma~\ref{lem:blocks}
is \texttt{block\_length} and \texttt{block\_zero}, also positional.

The converse computation of \texttt{sturmian.tex} \S6 is now legitimate. Suppose
$\ell$ is a coding whose blocks are $\seq{1}$ and $\seq{10}$, so $a_1 = 1$, and
suppose $\ell' = \ell$. By Corollary~\ref{cor:slope} the slopes agree, so
$\beta = \alpha$, and
\[
  \frac{1}{\alpha} \;=\; 1 + \alpha, \qquad \alpha^2 + \alpha - 1 = 0,
  \qquad \alpha = \frac{\sqrt 5 - 1}{2} .
\]

\section{Square roots}

Write $\alpha_0 = \alpha$ and $\alpha_{j+1} = \{1/\alpha_j\}$ for the tails of
the continued fraction, and $a_{j+1} = \lfloor 1/\alpha_j \rfloor$ for its
partial quotients. Let $\ell^{(j)}$ be the $j$th derived sequence of $\ell$.

\begin{theorem}\label{thm:tails}
Let $\ell = \C(\alpha, 0)$. For every $j \geq 1$,
\[
  \ell^{(j)} \;=\;
  \begin{cases}
    \Cp(\alpha_j,\, \alpha_j) & j \text{ odd} \\
    \C(\alpha_j,\, \alpha_j) & j \text{ even.}
  \end{cases}
\]
\end{theorem}

\begin{proof}
Let $\gamma$ be irrational in $(0,1)$, $a = \lfloor 1/\gamma\rfloor$ and
$\gamma' = \{1/\gamma\}$.

In $\C(\gamma,0)$ the first $1$ is at $k_0 = a$: for $k \leq a$ we have
$k\gamma \leq a\gamma < 1$ so $x_k = k\gamma$, and $k\gamma \geq 1-\gamma$ holds
first at $k = a$ because $(a+1)\gamma > 1$. Its coordinate is
\[
  u_0 \;=\; \frac{a\gamma - (1-\gamma)}{\gamma} \;=\; a + 1 - \frac{1}{\gamma}
  \;=\; 1 - \gamma' ,
\]
so $\{-u_0\} = \gamma'$ and Theorem~\ref{thm:derive} gives
$\C(\gamma,0)' = \Cp(\gamma',\gamma')$. With $\gamma = \alpha$ this is the case
$j = 1$.

In $\C(\gamma,\gamma)$ and in $\Cp(\gamma,\gamma)$ the first $1$ is at
$k_0 = a-1$, by the same count applied to $x_k = (k+1)\gamma$; the value
$x_{k_0} = a\gamma$ and the coordinate $u_0 = 1 - \gamma'$ are as above. So
Theorem~\ref{thm:derive} gives
\[
  \C(\gamma,\gamma)' = \Cp(\gamma',\gamma') , \qquad
  \Cp(\gamma,\gamma)' = \C(\gamma',\gamma') ,
\]
the second because $1 - u_0 = \gamma'$. Induction on $j$.
\end{proof}

\begin{corollary}\label{cor:cycle}
If the continued fraction of $\alpha$ is purely periodic with period $L$, then
$\ell = \C(\alpha,0)$ satisfies
\[
  \ell^{(L)} \;=\; \C(\alpha,\alpha) \;=\; \ell_1 \ell_2 \ell_3 \cdots ,
\]
the sequence $\ell$ with its first letter removed.
\end{corollary}

\begin{proof}
$\alpha_L = \alpha$, so Theorem~\ref{thm:tails} gives $\C(\alpha,\alpha)$ or
$\Cp(\alpha,\alpha)$, and those are equal by Lemma~\ref{lem:lowup} because the
orbit $\{(k+1)\alpha\}$ never meets $0$. Lemma~\ref{lem:shift} identifies
$\C(\alpha,\alpha)$ with the shift of $\ell$.
\end{proof}

\begin{corollary}\label{cor:cycleconv}
Conversely, if $\ell^{(L)}$ is a shift of $\ell$ then $\alpha_L = \alpha$, so
the continued fraction of $\alpha$ is purely periodic with period dividing $L$.
\end{corollary}

\begin{proof}
Shifting does not change the density, so Corollary~\ref{cor:slope} gives
$\alpha_L = \alpha$.
\end{proof}

Now take $n$ a non-square, $a_0 = \lfloor \sqrt n\rfloor$ and
$\alpha = \sqrt n - a_0$. If $\sqrt n = [\,a_0;\ \overline{a_1,\ldots,a_L}\,]$
then $\alpha = [\,0;\ \overline{a_1,\ldots,a_L}\,]$, purely periodic, since the
tails of $\alpha$ are the tails of $\sqrt n$ from index $1$. Pure periodicity is
Legendre's theorem, formalised in \texttt{lean/Palindrome.lean} as
\texttt{purely\_periodic}.

\begin{corollary}
For every non-square $n$, the sequence $\ell = \C(\sqrt n - a_0,\, 0)$ returns
to itself with one letter removed after exactly $L$ derivations, where $L$ is
the period of the continued fraction of $\sqrt n$.
\end{corollary}

\texttt{sturmian.tex} observes this for $n = 2$ and $n = 7$. For $n = 2$,
$\alpha = \sqrt 2 - 1$, $a_1 = 2$, $\beta = \alpha$ and $L = 1$: one derivation
returns the shift. For $n = 7$, $\alpha = \sqrt 7 - 2$ and the four tails are
\[
  \sqrt 7 - 2, \quad \frac{\sqrt 7 - 1}{3}, \quad \frac{\sqrt 7 - 1}{2},
  \quad \frac{\sqrt 7 - 2}{3}, \quad \sqrt 7 - 2 ,
\]
with partial quotients $1, 1, 1, 4$, which is the table of \S5 there. Every row
of that table is Sturmian by Corollary~\ref{cor:derivesturmian} and
Theorem~\ref{thm:complexity}. The intercept of each row equals its own slope,
by Theorem~\ref{thm:tails}.

\section{Lagrange}

\begin{theorem}
If the continued fraction of $\alpha$ is eventually periodic then $\alpha$ is a
quadratic irrational.
\end{theorem}

\begin{proof}
Let $\alpha_m = \alpha_{m+L}$. Each tail is a M\"obius image of the next,
$\alpha_j = 1/(a_{j+1} + \alpha_{j+1})$, so composing $L$ of them gives
integers $p,q,r,s$ with
\[
  \alpha_m \;=\; \frac{p\,\alpha_{m+L} + q}{r\,\alpha_{m+L} + s}
  \;=\; \frac{p\,\alpha_m + q}{r\,\alpha_m + s} ,
\]
hence $r\alpha_m^2 + (s-p)\alpha_m - q = 0$ with $(r,s-p,q) \neq 0$. So
$\alpha_m$ is a quadratic irrational, and composing $m$ further M\"obius maps
with integer coefficients keeps $\alpha$ in $\mathbb{Q}(\alpha_m)$.
\end{proof}

This is \texttt{lean/Lagrange.lean}, where it is stated for an abstract
sequence of tails, so it needs a witness: \texttt{Sturmian.sqrt\_quadratic} in
\texttt{lean/SqrtN.lean} supplies one. The Lean proves slightly more than the
display above. The leading coefficient $r$ is not merely part of a nonzero
triple but \emph{positive}, so the relation really has degree $2$; this comes
from the shape of the product, since
$\left(\begin{smallmatrix}p&q\\r&s\end{smallmatrix}\right)
\left(\begin{smallmatrix}0&1\\1&b\end{smallmatrix}\right) =
\left(\begin{smallmatrix}q&p+qb\\s&r+sb\end{smallmatrix}\right)$ --- the new
bottom-left entry is the old bottom-right one, so a positive bottom row stays
positive. For the transport to $\alpha$ itself the leading coefficient is the
form $(r,s,t)$ evaluated at the rational point $-S/R$, and that cannot vanish
because the roots of $(r,s,t)$ are $\alpha_m$ and its conjugate, neither of
which is rational.

\begin{theorem}[Lagrange, 1770]\label{thm:lagrange1770}
The continued fraction of every quadratic irrational is eventually periodic.
\end{theorem}

\begin{proof}
Write the complete quotients as $\alpha_k = (P_k + \sqrt D)/Q_k$ with
$Q_k \mid D - P_k^2$; the recurrence $P_{k+1} = a_kQ_k - P_k$,
$Q_{k+1} = (D - P_{k+1}^2)/Q_k$ preserves that divisibility, since
$P_{k+1} \equiv -P_k \pmod{Q_k}$. Call a state \emph{reduced} when
$\alpha_k > 1$ and $-1 < \bar\alpha_k < 0$. For a fixed $D$ there are only
finitely many reduced states, and on them the recurrence is invertible, so once
the orbit is reduced it is \emph{purely} periodic --- that is the argument of
\texttt{palindrome.tex} §2--3, which never used $Q = 1$. It is therefore enough
that some $\alpha_k$ be reduced.

For $k \geq 1$ we have $\alpha_k > 1$, so $a_k = \lfloor\alpha_k\rfloor \geq 1$.
Conjugating $\alpha_k = a_k + 1/\alpha_{k+1}$ gives
$\bar\alpha_k = a_k + 1/\bar\alpha_{k+1}$. Hence if $\bar\alpha_k < 0$ then
$1/\bar\alpha_{k+1} = \bar\alpha_k - a_k < -1$, so
$-1 < \bar\alpha_{k+1} < 0$ and $\alpha_{k+1}$ is reduced.

So suppose no $\alpha_k$ is reduced. Then $\bar\alpha_k > 0$ for every
$k \geq 1$ --- it is never $0$, as $\sqrt D$ is irrational. Positivity of
$\bar\alpha_{k+1} = 1/(\bar\alpha_k - a_k)$ forces $\bar\alpha_k > a_k \geq 1$,
so $\alpha_k$ \emph{and} $\bar\alpha_k$ exceed $1$ for every $k \geq 1$.
Subtracting the two conjugate recurrences, and using
$\alpha_k - a_k = 1/\alpha_{k+1}$ and
$\bar\alpha_k - a_k = 1/\bar\alpha_{k+1}$,
\[
  \alpha_{k+1} - \bar\alpha_{k+1}
  \;=\; \frac{1}{\alpha_k - a_k} - \frac{1}{\bar\alpha_k - a_k}
  \;=\; (\bar\alpha_k - \alpha_k)\,\alpha_{k+1}\bar\alpha_{k+1} .
\]
Since $\alpha_{k+1}\bar\alpha_{k+1} > 1$, the gap
$|\alpha_k - \bar\alpha_k|$ is strictly increasing. But
$\alpha_k - \bar\alpha_k = 2\sqrt D/Q_k$, so $|Q_{k+1}| < |Q_k|$ for every
$k \geq 1$: a strictly decreasing sequence of positive integers. That is
impossible.
\end{proof}

Two remarks. First, $\alpha_k\bar\alpha_k = -Q_{k-1}/Q_k$, so for $k \geq 1$
(where $\alpha_k > 0$) the condition $\bar\alpha_k < 0$ says exactly that
$Q_{k-1}$ and $Q_k$ have the \emph{same sign}. That is Lagrange's own
criterion, \emph{Additions} art. 41: "lorsque dans la série des nombres
$E, E_1, E_2,\ldots$ on en trouvera deux consécutifs de même signe, celui des
deux qui sera moindre que $\sqrt B$ sera déjà nécessairement périodique."

Second, the textbook proof is different: it passes through the convergents
$p_k/q_k$, the bound $|\alpha - p_k/q_k| < 1/q_k^2$, and either the invariance
of the discriminant of $A_kx^2 + B_kx + C_k$ or the uniqueness of the
continued-fraction expansion. The descent above needs none of that, which is
why the formalisation is short.

This is \texttt{lean/Quadratic.lean}: \texttt{exists\_reduced} is the preperiod,
\texttt{eventually\_periodic} the theorem. One thing the integer recurrence does
need there and did not for $\sqrt n$: $\lfloor (P + \sqrt D)/Q\rfloor$ equals
$\lfloor (P + a_0)/Q\rfloor$ only when $Q > 0$, and for a general quadratic
irrational the $Q_k$ change sign --- as the first remark says, they must. For
$Q < 0$ the value is $-\lfloor (P + a_0)/(-Q)\rfloor - 1$.

\section{What is not proved here}

Nothing. Both halves of Lagrange's theorem are proved and formalised. Euler's
is \S10, in \texttt{lean/Lagrange.lean}; Lagrange's own is
Theorem~\ref{thm:lagrange1770}, in \texttt{lean/Quadratic.lean}.

The numbers in the tables of \texttt{sturmian.tex} --- the slopes, the four
derived sequences of $n = 7$, the gap lists for $n = 2$ and $n = 7$, the four
factors of length $3$ --- have been checked in exact arithmetic in
$\mathbb{Q}(\sqrt n)$.

\end{document}
